\par\Needspace{6\baselineskip}\medskip
\hypertarget{native-097-article074}{}\pdfbookmark[2]{74.}{article-074}
\begin{center}74.\end{center}
Problema aliud, \emph{e corporis coelestis loco geocentrico atque situ plani orbitae eius locum heliocentricum in orbita derivare}, eatenus praecedenti affine est, quod quoque ab intersectione rectae inter terram et corpus coeleste ductae cum plano positione dato pendet. Solutio commodissime petitur e formulis art. 65, ubi characterum significatio haec erat:

$L$ longitudo terrae, $R$ distantia a Sole; latitudinem $B$ statuimus $=0$ (quum casus, ubi non est $=0$, ad hunc facile reduci possit per art. 72), unde $R'=R$; $l$ corporis coelestis longitudo geocentrica, $b$ latitudo, $\Delta$ distantia a terra, $r$ distantia a Sole, $u$ argumentum latitudinis, $\AscendingNode$ longitudo nodi ascendentis, $i$ inclinatio orbitae. Ita habemus aequationes
\[
\begin{aligned}
\mathrm{I.}\quad r\cos u-R\cos(L-\AscendingNode)&=\Delta\cos b\cos(l-\AscendingNode)\\[.5em]
\mathrm{II.}\quad r\cos i\sin u-R\sin(L-\AscendingNode)&=\Delta\cos b\sin(l-\AscendingNode)\\[.5em]
\mathrm{III.}\quad r\sin i\sin u&=\Delta\sin b.
\end{aligned}
\]
\noindent Multiplicando aequationem I per $\sin(L-\AscendingNode)\sin b$, II per $-\cos(L-\AscendingNode)\sin b$, III per $-\sin(L-l)\cos b$, fit additis productis
\[
\begin{aligned}
&\cos u\sin(L-\AscendingNode)\sin b\\
&\quad-\sin u\cos i\cos(L-\AscendingNode)\sin b\\
&\quad-\sin u\sin i\sin(L-l)\cos b=0,
\end{aligned}
\]
\noindent unde
\[
\mathrm{IV.}\quad \operatorname{tang}u=\frac{\sin(L-\AscendingNode)\sin b}{\cos i\cos(L-\AscendingNode)\sin b+\sin i\sin(L-l)\cos b}\mathpunct{\cdot}
\]
\noindent \SourceLeaf{098}Multiplicando autem I per $\sin(l-\AscendingNode)$, II per $-\cos(l-\AscendingNode)$, prodit productis additis
\[
\mathrm{V.}\quad r=\frac{R\sin(L-l)}{\sin u\cos i\cos(l-\AscendingNode)-\cos u\sin(l-\AscendingNode)}\mathpunct{\cdot}
\]
\noindent Ambiguitas in determinatione ipsius $u$ per aequ. IV sponte tollitur per aequ. III, quae ostendit, $u$ inter $0$ et $180^0$ vel inter $180^0$ et $360^0$ accipi debere, prout latitudo $b$ fuerit positiva vel negativa; sin vero fuerit $b=0$, aequatio V docet, statui debere $u=0$ vel $u=180^0$, prout $\sin(L-l)$ et $\sin(l-\AscendingNode)$ diversa signa habeant, vel eadem.

Computum numericum formularum IV et V variis modis per introductionem angulorum auxiliarium contrahere licet. E. g.
\[
\begin{aligned}
\text{statuendo}\quad &\frac{\operatorname{tang}b\cos(L-\AscendingNode)}{\sin(L-l)}=\operatorname{tang}A,\\[.4em]
\text{fit}\quad &\operatorname{tang}u=\frac{\sin A\operatorname{tang}(L-\AscendingNode)}{\sin(A+i)}
\end{aligned}
\]
\[
\begin{aligned}
\text{statuendo}\quad &\frac{\operatorname{tang}i\sin(L-l)}{\cos(L-\AscendingNode)}=\operatorname{tang}B,\\[.4em]
\text{fit}\quad &\operatorname{tang}u=\frac{\cos B\sin b\operatorname{tang}(L-\AscendingNode)}{\sin(B+b)\cos i}\mathpunct{\cdot}
\end{aligned}
\]
\noindent Perinde aequ. V per introductionem anguli cuius tangens $=\cos i\operatorname{tang}u$, vel $=\dfrac{\operatorname{tang}(l-\AscendingNode)}{\cos i}$ formam concinniorem nanciscitur. Sicuti formulam V e combinatione aequationum I, II obtinuimus, per combinationem aequationum II, III ad sequentem pervenimus:
\[
r=\frac{R\sin(L-\AscendingNode)}{\sin u(\cos i-\sin i\sin(l-\AscendingNode)\operatorname{cotang}b)}
\]
\noindent et perinde per combinationem aequationum I, III ad hanc:
\[
r=\frac{R\cos(L-\AscendingNode)}{\cos u-\sin u\sin i\cos(l-\AscendingNode)\operatorname{cotang}b}\mathpunct{\cdot}
\]
\noindent Utramque perinde ut V per introductionem angulorum auxiliarium simpliciorem reddere licet. Solutiones e praecedentibus demanantes collectae exemploque illustratae inveniuntur in \textsc{von Zach} \emph{Monatliche Correspondenz Vol. V. p. 540}\HistoricalNoteEight{Juni 1802. Band VI. S. 87.]}, quapropter hic evolutione ulteriori supersedemus. — Si praeter $u$ et $r$ etiam distantia $\Delta$ desideratur, per aequationem III determinari poterit.
